\section{Related Work and Scope of the Contribution}
\paragraph{Sensitivity-guided fault analysis.}
FKeras combines sensitivity analysis, fault injection, and selective protection using a projection derived from leading weight-Hessian eigenpairs~\cite{weng2024fkeras}. Its bit-tier-first ordering is an important baseline, not a diagonal-Hessian approximation: a full weight Hessian includes cross-parameter curvature. PrioriFI uses Hessian-informed initialization and measured degradation to prioritize subsequent injections~\cite{weng2026priorifi}. EISFINN assigns attribution-based sampling weights to neuron outputs and bits to find critical faults efficiently~\cite{colucci2024eisfinn}. These works establish that informed fault selection and bit prioritization are not new. Their discovery objectives do not automatically equal expected risk of a fixed protection mask under a prescribed simultaneous-fault distribution.

Progressive bit-flip attacks already exploit quantized weight encoding and gradient-guided, state-updating search~\cite{rakin2019bfa}. Their deliberately chosen damaging flips differ from independent Bernoulli events. Worst-case attack outcomes cannot be used as unweighted estimates of the latter expectation. Our experiments do not compare attack success rates or introduce an adversarial sampler.

\paragraph{Selective protection and contextual effects.}
Su et al. use DQN to choose layer-wise bit-position masks, with reward shaping
and transfer across layers~\cite{su2025rl}. Each selected position protects that
position in every weight of a layer. Their simulated memory-entry fault model
and performance/cost reward differ from coordinate-wise Bernoulli faults and a
fixed global bit budget here. This is a learned allocation policy, not an
isolated sensitivity score. Its interaction-dependent feedback can represent
multi-fault behavior; our analytical result does not rule that out. A matched
reimplementation would require common fault laws, granularity and training cost;
we do not insert unmatched published numbers into our curves.
Bit-level redundancy work exploits structure in stored representations~\cite{catalan2025redundancy}, while the public \emph{Learned or Ranked?} companion release compares learned and rule-based allocation, search methods, and layer-interaction diagnostics~\cite{alobaidi2026selectiveeccdataset}. We cite the latter as a dataset/software release; peer-reviewed acceptance of its associated manuscript is not established here. Its contextual allocation analysis is relevant prior art regardless of that status. Its layer granularity, codeword protection semantics, and fault generators differ from our exact coordinate-Bernoulli formulation. We claim neither the first observation of nonadditivity nor the first context-dependent protection decision.

\paragraph{Extension beyond the earlier NetSuRF flow.}
The FCCM predecessor~\cite{guo2025inpixel} described bit-level susceptibility and a gradient/delta technique within a broader selective-protection flow. The added contribution here is not a new physical flow: it is the explicit finite-code contract and experimental separation of prediction quality from allocation quality across three tasks. The present models and measurements are not reproductions of that paper's unspecified checkpoint or its reported implementation results. Historical qPolarGrad is separately identified in Section~\ref{sec:estimators}; no old superiority claim is inherited.

\paragraph{What is, and is not, new.}
Multilinear expectations and conditional importance belong to established Boolean-function and reliability theory~\cite{odonnell2014analysis,rauzy2022reliability}. We use a simple specialization to explain a decision failure and to define bounded diagnostic comparisons. The study's distinction is the controlled connection between exact encoded perturbations, familiar ranking rules, and held-out task outcomes at equal logical budgets. It is not a claim to invent bit injection, gradient sensitivity, importance sampling, or generic interaction analysis. Nor do the results establish comparative efficacy against physical error-correction or redundancy implementations; those are outside the software-only experiment.
